\begin{lemma}
\label{editorial-source-lemma-same-image}
Let $k$ be a field, and let $R$, $S$ be $k$-algebras.
Let $S' \subset S$ be a sub $k$-algebra, and let $f \in S' \otimes_k R$.
In the commutative diagram
$$
\xymatrix{
\Spec((S \otimes_k R)_f) \ar[rd] \ar[rr] & &
\Spec((S' \otimes_k R)_f) \ar[ld] \\
& \Spec(R) &
}
$$
the images of the diagonal arrows are the same.
\end{lemma}

\begin{proof}
Retain the original inclusion $i:S'\hookrightarrow S$,
the original $k$-algebra $R$, and the specified
$f\in S'\otimes_kR$. For a prime $\mathfrak p\subset R$
put $K=\kappa(\mathfrak p)$. The coefficient maps
$R\to K$ give elements $f_{S'}\in S'\otimes_kK$
and $f_S\in S\otimes_kK$, related by $i\otimes1$.
For either $B=S'$ or $B=S$, the exact comparison is
$$
((B\otimes_kR)_f)\otimes_RK
\longrightarrow (B\otimes_kK)_{f_B},
\qquad
\frac{\sum_j b_j\otimes r_j}{f^d}\otimes\lambda
\longmapsto
\frac{\sum_j b_j\otimes\overline{r_j}\lambda}{f_B^d}.
$$
Here the $f$ in the $B=S$ expression denotes its
image under the original inclusion, not a different
choice of element. The inverse sends
$$
\frac{b\otimes\lambda}{f_B^d}
\longmapsto \frac{b\otimes1}{f^d}\otimes\lambda
$$
and extends additively. To check these maps, first use
the balanced isomorphism
$(B\otimes_kR)\otimes_RK\simeq B\otimes_kK$,
$(b\otimes r)\otimes\lambda\mapsto b\otimes\overline r\lambda$,
whose inverse is $b\otimes\lambda\mapsto(b\otimes1)\otimes\lambda$.
On localizing, $f$ maps to $f_B$ and is inverted in
both targets. The universal property of localization
therefore extends the two inverse algebra maps to
the displayed fraction maps. In particular it checks
all zero-divisor witnesses as well as tensor balancing;
their composites are identities on the original
algebra generators and on $1/f$.

For a ring $C$ and $h\in C$, one has $C_h=0$
exactly when $h$ is nilpotent. Indeed $1/1=0$
holds exactly when $h^d\cdot1=0$ for some
$d\geq0$; in the case $d=0$ the ring is zero
and $h^1=0$ also holds. Conversely any such positive
power makes $h$ an invertible nilpotent, forcing the
localized unit to vanish. Thus the fibre criterion
of Lemma \ref{editorial-source-lemma-in-image} identifies the image
of each diagonal arrow with the primes for which
its original $f_B$ is not nilpotent.

Since $k$ is a field, tensoring the injection
$S'\hookrightarrow S$ with the $k$-vector space $K$
is injective. Explicitly, choose a $k$-basis of $K$.
Under the direct-sum identifications of both tensor
modules, the map is the direct sum of copies of the
original injection; hence its kernel is zero.
Consequently
$$
f_{S'}^d=0\quad\Longleftrightarrow\quad f_S^d=0
\qquad(d\geq1).
$$
The forward implication is functoriality; the reverse
uses this injection. Nonnilpotence is therefore the
same in both original fibres. The fibre criterion
and the fraction comparisons prove equality of the
two spectrum images. Independently, one inclusion
is also the factorization through the horizontal
arrow: every prime of the larger tensor algebra
contracts to a prime of the smaller one, with the
same contraction to $R$. The reverse inclusion is
the nonnilpotence argument just supplied.

\medskip\noindent
This proof gives a more general assertion with its
exact sufficient hypothesis. Let $A$ be any base
ring, let $B\to C$ be a map of $A$-algebras, and
let $R$ be an $A$-algebra. Suppose that for every
original $\mathfrak p\in\Spec(R)$ the map
$$
B\otimes_A\kappa(\mathfrak p)
\longrightarrow C\otimes_A\kappa(\mathfrak p)
$$
is injective. Then for every $f\in B\otimes_AR$,
the two original principal-localization spectrum
images in $\Spec(R)$ agree. All the preceding
fraction maps work with $A$ in place of $k$,
and injectivity of these particular maps supplies
exactly the nilpotence reflection used in the proof.
In particular the hypothesis holds when the
$A$-module map $B\to C$ is pure, meaning that
$B\otimes_AL\to C\otimes_AL$ is injective for
every $A$-module $L$. It suffices, more weakly,
that each of the displayed fibre maps reflect
nilpotence of the particular image of $f$.

An arbitrary injection over a general base ring
does not suffice. Take $A=R=B=\mathbf Z$,
$C=\mathbf Q$, and $f=1$. The spectrum image
for $B$ is all of $\Spec(\mathbf Z)$, whereas
the image for $C$ is $\{(0)\}$. At a nonzero
prime $(p)$ the receiving map of fibres is
$\mathbf F_p\to0$, since
$q\otimes\overline a=(q/p)\otimes p\overline a=0$.
This identifies the failed residue injection and
the failed image equality for the same original
maps and element. Finally, nonzero elements cannot
replace nonnilpotent elements in the proof: over
$k[\epsilon]/(\epsilon^2)$ the nonzero element
$\epsilon$ has empty principal open. Its square
is zero and its localization is the zero ring.
\end{proof}